Foundation model-based multi-agent systems now execute production workloads whose reliability depends on two layers that the research community has built separately: diagnostics, which explains failure after the fact, and infrastructure, which keeps serving efficient in the moment. This review has assembled both into a single taxonomy: four diagnostic subcategories spanning three failure-localization paradigms and three safety paradigms, and four infrastructure optimization axes spanning memory optimization, prefix-aware scheduling, cache lifecycle, and power management. It has compared the reviewed systems quantitatively where the evidence permits, holding publication type and verification tier explicit throughout.

The central argument is that these layers are not independent. They act on the same runtime objects under opposed objectives: diagnostics must retain state long enough to explain a failure, while serving infrastructure must discard it as soon as reuse becomes unlikely. We have shown that this tension is already visible in two cross-layer mechanisms and an intra-layer analogue that hold for post-hoc diagnosis, with a boundary condition marking where it would become acute online, and that no surveyed system accounts for retention as a managed resource. Treating retention as a shared, budgeted resource, rather than a local eviction heuristic, is we argue the problem that connects the two layers and the most promising direction for the work that follows. The taxonomy, the verified comparisons, and the tension analysis presented here are offered as a common frame for that effort.
